\documentclass[11pt,a4paper]{article}
\usepackage{gtu-el-yazisi}
\newcommand\daire[1]{\tikz[baseline=(d.base)]\node[draw=murekkep,circle,inner sep=0.2pt](d){$#1$};}
\begin{document}
\notbilgi{Diferansiyel Denklemler}{1. Hafta -- Diferansiyel denklemlere giriş (teori)}{diferansiyel-teori-1-hafta-notum.pdf}

\section{Diferansiyel denklemler}
\sayfa{1}
\begin{center}
\begin{tabular}{cc}
$\swarrow$ & $\searrow$\\
ODE & PDE
\end{tabular}
\end{center}
\begin{tanim}
$x$ serbest değişken, $y=y(x)$ ona bağlı bir fonksiyon ve bu fonksiyonun $y',y'',\cdots y^{(n)}$ arasındaki
\[ F(x,y,y',y'',\cdots y^{(n)}) = 0 \qquad (1) \]
şeklindeki denklemlere adi diferansiyel denklem denir.

eğer $F$ fonksiyonu $y$ ve onun türevlerine göre yani $(y',y'',\cdots,y^{(n)})$ lineer olursa bu durumda $(1)$'e lineer adi dif denklem denir
\end{tanim}

Örneğin: 1. mertebeden lineer adi diferansiyel denklem en genel şekli
\[ a(x)\,\daire{y'} + b(x)\,\daire{y} = c(x) \]
2. mertebeden lineer adi dif den en genel şekli
\[ a_1(x)\,y''(x) + a_2(x)\,y' + a_3(x)\,y = b(x) \quad \text{olur} \]

\begin{ornek}
\begin{tabular}{@{}l@{\qquad}l@{}}
$3y\,y' + y^2 - 2x$ & lineer olmayan birinci mertebeden adi dif denk\\[1mm]
$(1+x^2)\,\daire{y'} + 2x\,\daire{y} = 3$ & lineer birinci mertebeden adi dif denk.\\[1mm]
$\daire{y''} + 2x\,\daire{y'} + x^2\,\daire{y} = \sin x$ & 2. mertebeden lineer adi dif denk\\[1mm]
$y''' + 3y^{\daire{2}} = 5$ & 3. \ \ ,, \ \ lineer olmayan adi dif denklem
\end{tabular}
\end{ornek}

$(1)$ denkleminde $n=1$ alırsak
\[ F(x,y,y') = 0 \]
1. mertebeden dif denk en genel halde elde edilir.

Eğer bu denklem $y$'nin türevine göre çözülebiliyorsa $y'=f(x,y)$ şeklinde yazılabiliyorsa 1. mertebeden açık adi dif denk denir.

\sayfa{2}
Eğer $y=\varphi(x)$ fonksiyonu $I=(\alpha,\beta)$ aralığında $n$. mertebeden sürekli türevlere sahip ve her $x\in I$ için
\[ F(x,\varphi(x),\varphi'(x),\cdots,\varphi^{(n)}(x)) = 0 \]
sağlıyorsa, bu durumda $y=\varphi(x)$ fonksiyonu $I$ aralığında $(1)$ denkleminin çözümüdür.

Örneğin: \ $y = x + \dfrac{1}{\sqrt{2}}$ fonksiyonu $I$ aralığında
\[ y = xy' + \frac{y'}{\sqrt{1+(y')^2}} \qquad \text{denkleminin çözümüdür.} \]
\[ y' = 1 \qquad\qquad x + \frac{1}{\sqrt{2}} = x + \frac{1}{\sqrt{2}} \]

\begin{tanim}
$(1)$ denklemini $\varphi(x,y)=0$ şeklindeki kapalı fonksiyon çözümüne onun integrali denir.
\end{tanim}

En basit adi dif denklemler;
\begin{enumerate}[label=\protect\daire{\arabic*}]
\item $y' = f(x)$ \quad (veya $\dfrac{dy}{dx} = f(x)$) \hfill {\small normalde $y=f(x,y)$}

Eğer $f(x)$ fonksiyonunun ilkelleri $F(x)+c$, $c=$ sabit ise o halde bu denklemin çözümü
\[ y = \int f(x)\,dx + C = F(x) + C \]
{\small Örnek: \ $y'=2x$, \ $y=\int 2x\,dx = x^2+C$}

\textbf{Not:} bir adi dif denklemin sonsuz sayıda çözümü vardır.

\item $y' = f(y)$ \quad $\left(\dfrac{dy}{dx} = f(y)\right)$
\[ \frac{dy}{dx} = f(y) \;\Rightarrow\; \frac{dx}{dy} = \frac{1}{f(y)}, \qquad f(y)\neq 0 \]
\[ x = \underbrace{\int \frac{1}{f(y)}\,dy + C}_{G(y)+C} \]
Burada $G(y)$ fonksiyonu $\dfrac{1}{f(y)}$'nin ilkelidir
\end{enumerate}

\sayfa{3}
\begin{ornek}
$y' = x + \sin x$
\[ y = \frac{x^2}{2} - \cos x + C \]
\end{ornek}
\begin{ornek}
$y' = y^2 + y$
\end{ornek}
\[ \frac{dy}{dx} = y^2 + y \]
\[ \Rightarrow\; \frac{dx}{dy} = \frac{1}{y^2+y} \qquad\qquad y^2+y\neq 0 \]
\begin{align*}
x &= \int \frac{1}{y^2+y}\,dy + C\\
&= \int \left(\frac{1}{y} - \frac{1}{y+1}\right)dy + C\\
&= \ln|y| - \ln|y+1| + C
\end{align*}
\[ \frac{1}{y^2+y} = \frac{1}{y(y+1)} = \frac{\overset{=1}{A}}{y} + \frac{\overset{=-1}{B}}{y+1} \]
\[ x = \ln\left|\frac{y}{y+1}\right| + C, \qquad C\in\mathbb{R} \]
\[ y=0,\ y=-1 \]
\[ y^2 + y = 0 \;\Rightarrow\; y=0,\ y=-1 \]
\begin{center}
denklemde koyarız\\
denklemi sağladıkları için çözümdür
\end{center}

\subsection{\texorpdfstring{\protect\daire{1.2}}{1.2} Cauchy Problemi (veya başlangıç değer problemi)}
\noindent\begin{minipage}[t]{0.68\linewidth}
$y = f(x,y)$ denklemini alalım. $f(x,y)$ fonksiyonu bir $D\subset\mathbb{R}^2$ bölgesinde tanımlı olsun. Bir $(x_0,y_0)\in D$ noktasını alalım, $y'=f(x,y)$ denkleminin $y(x_0)=y_0$ koşulunu sağlayan çözümünün bulunması problemine Cauchy problemi / başlangıç değer problemi denir.
\end{minipage}\hfill
\begin{minipage}[t]{0.28\linewidth}\centering\vspace{0pt}
\begin{tikzpicture}[scale=0.8,murekkep]
\draw[->] (-0.4,0) -- (2.2,0);
\draw[->] (0,-0.9) -- (0,1.4);
\draw (0.9,0.2) circle (0.95);
\fill (0.6,0.45) circle (1.3pt) node[above right=-1pt] {\scriptsize $(x_0,y_0)$};
\node at (1.35,-0.35) {\small $D$};
\node at (2.1,1.2) {\small $\mathbb{R}^2$};
\end{tikzpicture}
\end{minipage}
\[ \left.\begin{aligned} y' &= f(x,y) \ \cdots & &(2a)\\ y(x_0) &= y_0 & &(2b) \end{aligned}\right\}\ (2) \]
$x_0$, $y_0$ değerlerine başlangıç verileri, $(2b)$'ye başlangıç koşulu denir.

\sayfa{4}
Şimdi: \ $y'=f(x,y)$, \ $\dfrac{dy}{dx}=f(x,y)$ \ denklemi için bazı çözüm yöntemleri belirleyelim
\[ f(x,y) = -\frac{M(x,y)}{N(x,y)} \quad \text{olduğunu kabul edersek denklemi simetrik şekilde} \]
\[ M(x,y)\,dx + N(x,y)\,dy = 0 \]
denk denklemine dönüşür

\subsection{\texorpdfstring{\protect\daire{2}}{2} Ayrılabilir denklemler}
\[ y' = f_1(x)\,f_2(y) \quad \text{şeklinde veya} \]
\[ M_1(x)\,M_2(y)\,dx + N_2(y)\cdot N_1(x)\,dy = 0 \quad \text{şeklindeki denklemlere} \]
ayrılabilir denklemler denir.

Eğer $M_2(y)\,N_1(x)\neq 0$ ise
\[ \frac{M_1(x)\,M_2(y)\,dx}{M_2(y)\,N_1(x)} + \frac{N_2(y)\,N_1(x)\,dy}{M_2(y)\,N_1(x)} = 0 \]
\[ \frac{M_1(x)}{N_1(x)}\,dx + \frac{N_2(y)}{M_2(y)}\,dy = 0 \]
elde edilir. Bu eşitliği integrallersek
\[ \int \frac{M_1(x)}{N_1(x)}\,dx + \int \frac{N_2(y)}{M_2(y)}\,dy = C, \qquad \forall C\in\mathbb{R} \]
elde edilir ki bu da ayrılabilir dif denk çözümüdür

$M_2(y)=0$, $N_1(x)=0$ olduğu durumlara ayrıca kontrol edilmelidir

\begin{ornek}
$x(1+y^2) + y(1+x^2)\dfrac{dy}{dx} = 0$ denkleminin integralleyiniz
\end{ornek}
\begin{gather*}
x(1+y^2)\,dx + y(1+x^2)\,dy = 0\\
(1+y^2)(1+x^2)\neq 0\\
\frac{x}{1+x^2}\,dx + \frac{y}{1+y^2}\,dy = 0\\
\int \frac{x}{1+x^2}\,dx + \int \frac{y}{1+y^2}\,dy = 0\\
\frac{1}{2}\int \frac{2x}{1+x^2}\,dx + \frac{1}{2}\int \frac{2y}{1+y^2}\,dy = 0\\
\frac{1}{2}\ln(1+x^2) + \frac{1}{2}\ln(1+y^2) = c
\end{gather*}

\sayfa{5}
\[ \ln(1+x^2)(1+y^2) = 2C \]
\[ \underbrace{(1+x^2)(1+y^2) = c_1}_{\text{çözüm ailesi}}, \qquad c_1 = e^{2C} > 0 \]

\begin{ornek}
$\begin{cases} y' = xy^2 + 2xy\\ y(0) = 1 \end{cases}$ \qquad başlangıç değer problemini çözünüz.
\end{ornek}
\begin{gather*}
\frac{dy}{dx} = x(y^2+2y)\\
dy - x(y^2+2y)\,dx = 0\\
y^2+2y\neq 0,\\
\frac{dy}{y^2+2y} - x\,dx = 0\\
\int \frac{dy}{y^2+2y} - \int x\,dx = C, \qquad C\in\mathbb{R}\\
\frac{1}{y^2+2y} = \frac{1}{y(y+2)} = \frac{\overset{=\frac12}{A}}{y} + \frac{\overset{=-\frac12}{B}}{y+2}\\
\frac{1}{2}\int\left(\frac{1}{y} - \frac{1}{y+2}\right)dy - \frac{x^2}{2} = 0\\
\frac{1}{2}\bigl(\ln|y| - \ln|y+2|\bigr) - \frac{x^2}{2} = c\\
\ln\left|\frac{y}{y+2}\right| - x^2 = 2C \qquad\qquad \left|\frac{y}{y+2}\right| = e^{x^2}e^{2C}\\
\frac{y}{y+2} = \mp e^{2C_1}e^{x^2}
\end{gather*}
\[ \Rightarrow\ \frac{y}{y+2} = Ce^{x^2}, \quad C=\mp e^{2C_1}
\qquad\longrightarrow\qquad
\begin{aligned} y &= yCe^{x^2} + 2Ce^{x^2}\\ y(1-Ce^{x^2}) &= 2Ce^{x^2} \end{aligned} \]
\[ \Rightarrow\ y = \frac{2Ce^{x^2}}{1-Ce^{x^2}}, \qquad \text{burada } C=\mp e^{2C_1}\neq 0 \]
\[ y^2+2y = 0 \;\rightarrow\; y(y+2) = 0 \;\rightarrow\;
\begin{array}{l} y=0\ \checkmark \ \text{çözüm}\\ y=-2\ \checkmark\ \Bigl(\begin{array}{l}0 = 4x-4x\\ 0=0\end{array}\Bigr)\ \checkmark \end{array} \]

\sayfa{6}
\[ y(\overset{x_0}{0}) = \overset{y_0}{1} \]
\begin{gather*}
1 = \frac{2Ce^0}{1-Ce^0}\\
1 - C = 2C\\
3C = 1\\
C = \frac{1}{3}
\end{gather*}
Başlangıç değer probleminin çözümü
\[ y = \frac{2\cdot\frac{1}{3}e^{x^2}}{1-\frac{1}{3}e^{x^2}} = \frac{2e^{x^2}}{3-e^{x^2}} \quad \text{olur} \]

\begin{ornek}
$e^x\,dx - (1+e^x)\,y\,dy = 0$ denkleminin $(0,1)$ noktasından geçen integral eğrisini bulunuz
\end{ornek}
\[ \begin{cases} e^x\,dx - (1+e^x)\,y\,dy = 0\\ y(1) = 1\footnotemark \end{cases} \]
\footnotetext{Claude notu: $(0,1)$ noktası için muhtemelen $y(0)=1$ olmalı (aşağıda $x=0$ kullanılmış).}
\begin{gather*}
\frac{e^x}{1+e^x}\,dx - y\,dy = 0\\
\int \frac{e^x}{1+e^x}\,dx - \int y\,dy = C\\
\ln(1+e^x) - \frac{y^2}{2} = 0\\
y^2 = 2\ln(1+e^x) - 2C\\
y^2 = 2\ln(1+e^x) + C, \qquad \forall C\in\mathbb{R}\\
1 = 2\ln(1+e^0) + C\\
C = 1 - 2\ln(2)\\
\Rightarrow\ y^2 = 2\ln(1+e^x) + 1 - 2\ln(2) = 2\ln\left(\frac{1+e^x}{2}\right) + 1
\end{gather*}

\subsection{\texorpdfstring{\protect\daire{3}}{3} Ayrılabilen denkleme indirgenebilen denklemler}
\sayfa{7}
\daire{3.1} \quad $y' = f(ax+by+c)$

denklemi $ax+by+c = z$, \ $z(x)$ dönüşümüyle $z' = a + by' \Rightarrow y' = \dfrac{z'-a}{b}$ olduğuna göre
\[ z' = a + b\,f(z) \]
şeklinde dönüşür, bu bir ayrılabilir denklemdir

\begin{ornek}
$y' = \cos(x-y-1)$ denklemini çözünüz
\end{ornek}
\begin{cozum}
$z = x-y-1$, \ $z=z(x)$
\[ z' = 1 - y' \;\rightarrow\; y' = 1 - z' \]
yerine yazarsak,
\begin{gather*}
1 - z' = \cos(z)\\
z' = 1 - \cos z\\
\frac{dz}{dx} = 1 - \cos z \;\Rightarrow\; dz - (1-\cos z)\,dx = 0\\
1-\cos z\neq 0 \ \text{ise} \ \int \frac{dz}{1-\cos z} = \int dx
\end{gather*}
\[ \kutula{\cos 2z = 1 - 2\sin^2 z} \qquad \cos z = 1 - 2\sin^2\frac{z}{2} \qquad 2\sin^2\frac{z}{2} = 1 - \cos z \]
\[ \int \frac{dz}{2\sin^2\frac{z}{2}} = x + C \;\Rightarrow\; -\cot\left(\frac{z}{2}\right) = x + c \]
\[ \cot\left(\frac{x-y-1}{2}\right) = -x - c \]
\stepcounter{footnote}%
\begin{gather*}
\frac{1}{2}\int \frac{dz}{\sin^2\frac{z}{2}} = \frac{1}{2}\int \frac{du}{\sin^2 u}\,{}^{\thefootnote}\\
= -\cot u\\
= -\cot\left(\frac{z}{2}\right)
\end{gather*}
\[ 1 - \cos z = 0 \Rightarrow \cos z = 1 \Rightarrow z = 2n\pi, \quad n=0,1,2,\ldots \]
\[ x - y - 1 = 2n\pi \Rightarrow y = x - 2n\pi - 1 \quad \text{böylece çözüm ailesi} \]
\[ \cot\left(\frac{x-y-1}{2}\right) = -x + c, \quad \forall c\in\mathbb{R} \]
\[ y = x - 2n\pi - 1 \qquad n = 0,1,2,\ldots \]
\end{cozum}
\footnotetext{Claude notu: $u=z/2$, $dz=2\,du$ olduğundan muhtemelen $\frac12\int\frac{dz}{\sin^2(z/2)}=\int\frac{du}{\sin^2u}$ olmalı (sonuç $-\cot u$ bununla uyumlu).}

\sayfa{8}
\daire{3} Ayrılabilen denkleme indirgenebilen denklemler \hfill $\begin{array}{r} f(tx,ty) = f(x,y)\\ \forall t\in\mathbb{R} \end{array}$

\subsection{\texorpdfstring{\protect\daire{4}}{4} Homojen denklemler;}
$y' = f\left(\dfrac{x}{y}\right)$\footnote{Claude notu: aşağıdaki $\frac{y}{x}=u$ dönüşümüne göre muhtemelen $y'=f\left(\frac{y}{x}\right)$ olmalı.} şeklindeki denklemlere homojen denklemler denir.

bu denklem $\dfrac{y}{x} = u$, \ $u(x) = u$, dönüşümü ile ayrılabilir denk dönüştürülür. $\dfrac{y}{x} = u \Rightarrow y = xu \Rightarrow y' = u + x\cdot u'$

yerine yazarsak \quad $u + xu' = f(u)$
\[ x\,\frac{du}{dx} = f(u) - u \]
\[ f(u) - u \neq 0 \quad \text{------} \quad \frac{du}{f(u)-u} = \frac{dx}{x} \]
ve bu eşitliği integrallersek çözüm elde ederiz

\end{document}
